\documentclass[a4paper,11pt]{ctexart}

\usepackage[top=1.8cm,bottom=1.8cm,left=2.1cm,right=2.1cm]{geometry}
\usepackage{amsmath,amssymb}
\usepackage{enumitem}
\usepackage{fancyhdr}
\usepackage{lastpage}

\pagestyle{fancy}
\fancyhf{}
\renewcommand{\headrulewidth}{0pt}
\fancyfoot[C]{\small 第~\thepage~页\quad 共~\pageref{LastPage}~页}

\linespread{1.28}
\setlength{\parindent}{0pt}

\newcommand{\tr}{\operatorname{tr}}
\newcommand{\rk}{\mathrm{r}}
\newcommand{\R}{\mathrm{R}}
\newcommand{\T}{\mathrm{T}}
\newcommand{\blank}[1][2.8cm]{\underline{\hspace{#1}}}

\begin{document}

\begin{center}
  {\Large\bfseries 华中科技大学《线性代数》}\\[0.4em]
  {\Large\bfseries 2018--2019 学年第二学期期末考试 A 卷}
\end{center}

\vspace{0.8em}

\textbf{一、判断题（每小题 2 分，共 16 分）}

\begin{enumerate}[label=\arabic*., leftmargin=2em, itemsep=0.55em, topsep=0.5em]
  \item 设 $A, B, C$ 是三个 $n$ 阶非零矩阵，如果 $AB = BA$，则 $A^2 B + ACA = A(B + C)A$.

  \item 若向量 $\alpha_1, \alpha_2, \dots, \alpha_m$ 线性无关，则必有 $\alpha_1 + \alpha_2 + \dots + \alpha_m \neq 0$.

  \item 设 $A, B$ 为满足 $AB = 0$ 的两个非零矩阵，则 $A$ 的列向量组线性相关，$B$ 的列向量组也线性相关.

  \item 设 $A$ 为三阶非零矩阵满足 $A^2 = 0$，则齐次线性方程组 $AX = 0$ 的线性无关的解向量的个数为 $2$.

  \item 设 $A$ 为 $m \times n$ 阶实矩阵，则矩阵 $A^{\T}A$ 与 $A$ 有相同的秩，其中 $A^{\T}$ 是 $A$ 的转置.

  \item 若方阵 $A$ 的各行元素和为 $1$，则 $1$ 必为 $A$ 的一个特征值.

  \item 记 $A^*$ 为方阵 $A$ 的伴随矩阵，若 $A$ 为正定矩阵，则 $A^*$ 也为正定矩阵.

  \item 设 $A$ 是 $n$ 阶实反对称矩阵，即 $A^{\T} = -A$，其中 $A^{\T}$ 是 $A$ 的转置，则对任意的 $n$ 维实向量 $X$，均有 $X^{\T}AX = 0$.
\end{enumerate}

\vspace{0.9em}

\textbf{二、填空题（每小题 4 分，共 20 分）}

\begin{enumerate}[label=\arabic*., leftmargin=2em, itemsep=0.9em, topsep=0.5em]
  \item 设 $A$ 是 $n$ 阶方阵，$I$ 为 $n$ 阶单位矩阵，且 $A^2 = A$，则 $(A - 2I)^{-1} = \blank$.

  \item 设 $A$ 是 $n$ 阶可逆矩阵，$A$ 的行列式 $|A| = a$，且 $A$ 的各行元素之和均为 $b$，则 $A$ 的第 $1$ 列元素的代数余子式之和 $A_{11} + A_{21} + \dots + A_{n1} = \blank$.

  \item 设 $\alpha_1, \alpha_2, \alpha_3$ 为三个三维（列）向量，记矩阵 $A = (\alpha_1, \alpha_2, \alpha_3)$，已知 $A$ 的行列式 $|A| = 1$，\\
  记矩阵 $B = (\alpha_1 + \alpha_2 + \alpha_3,\ \alpha_1 + 2\alpha_2 + 4\alpha_3,\ \alpha_1 + 3\alpha_2 + 9\alpha_3)$，则行列式 $|B| = \blank$.

  \item 已知 $f(x_1, x_2, x_3) = x_1^2 + t x_2^2 + t x_3^2 + 2 x_1 x_2 + 4 x_1 x_3$ 为正定二次型，\\
  则 $t$ 的取值范围为 $\blank$.

  \item 设矩阵 $A = (a_{ij})_{3 \times 3}$ 是正交矩阵，且 $a_{11} = 1$. 设向量 $\beta = (1, 0, 0)^{\T}$，则非齐次线性方程组 $AX = \beta$ 的解 $X = \blank$.
\end{enumerate}

\vspace{0.9em}

\textbf{三、（8 分）}计算 $n$ 阶行列式
\[
  D_n = \begin{vmatrix}
    1 + a_1 & 1 & \cdots & 1 \\
    1 & 1 + a_2 & \cdots & 1 \\
    \vdots & \vdots & \ddots & \vdots \\
    1 & 1 & \cdots & 1 + a_n
  \end{vmatrix},
  \quad\text{其中 } a_i \neq 0,\ i = 1, 2, \dots, n.
\]

\vspace{0.7em}

\textbf{四、（12 分）}设线性方程组为
\[
  \begin{cases}
    x_1 + x_2 - 2x_3 + 3x_4 = 0, \\
    2x_1 + x_2 - 6x_3 + 4x_4 = -1, \\
    3x_1 + 2x_2 + a x_3 + 7x_4 = -1, \\
    x_1 - x_2 - 6x_3 - x_4 = b.
  \end{cases}
\]
\begin{enumerate}[label=(\arabic*), leftmargin=2.4em, itemsep=0.25em, topsep=0.3em]
  \item 讨论 $a, b$ 为何值时，方程组有解或无解；
  \item 当有解时，求出方程组的解.
\end{enumerate}

\vspace{0.7em}

\textbf{五、（10 分）}设矩阵 $A, B$ 满足 $A^* B A = 2BA - 8I$，其中
\[
  A = \begin{pmatrix}
    1 & 2 & -2 \\
    0 & -2 & 4 \\
    0 & 0 & 1
  \end{pmatrix},
\]
$A^*$ 表示 $A$ 的伴随矩阵，$I$ 为三阶单位矩阵，求矩阵 $B$.

\vspace{0.7em}

\textbf{六、（12 分）}已知矩阵
\[
  A = \begin{pmatrix}
    a - 1 & -1 & 1 \\
    -1 & a - 1 & 1 \\
    1 & 1 & a - 1
  \end{pmatrix}
\]
有特征值 $1$，当 $A$ 的特征值之和最小时，求正交矩阵 $C$，使得 $C^{\T} A C$ 为对角矩阵，并写出此对角矩阵.

\vspace{0.7em}

\textbf{七、（10 分）}已知方阵 $A,\ I - A,\ I - A^{-1}$ 均可逆，其中 $I$ 为单位矩阵，证明：
\[
  (I - A)^{-1} + (I - A^{-1})^{-1} = I.
\]

\vspace{0.7em}

\textbf{八、（12 分）证明题}

设 $A, B$ 为 $n$ 阶方阵，已知 $A$ 可相似对角化，又有 $A^2 + A = 0$，$B^2 + B = I$，且有秩关系 $\R(AB) = 2$.
证明：
\begin{enumerate}[label=(\arabic*), leftmargin=2.4em, itemsep=0.25em, topsep=0.3em]
  \item $B$ 可逆；
  \item $|A + 2I| = 2^{n-2}$，其中 $I$ 为 $n$ 阶单位矩阵.
\end{enumerate}

\end{document}
